In this section we develop the properties of the tensor Casimir element that are needed in the $r$-matrix formalism \cite{perelomov1990integrable,babelon2003introduction}.

The $r$-matrix formalism is naturally expressed in terms of tensorial objects associated with Lie algebras. Among these, a distinguished role is played by the Casimir element. It provides a canonical invariant tensor and will be used to constrain the structure of $r$-matrices.

We now recall its definition and main properties.

\begin{definition}
    The Casimir tensor element of $\mathfrak{gl}(n)$ associated with the trace form is defined as
    \begin{equation}\label{eq:casimir_element}
        C_{12} = \sum_{i,j} E_{ij} \otimes E_{ji}.
    \end{equation}
\end{definition}
\begin{proposition}\label{prop:casimir_invariant}
    The Casimir tensor element satisfies
    \begin{equation}\label{eq:casimir_property}
        [C_{12}, X \otimes 1 + 1 \otimes X] = 0
        \qquad \forall X \in \mathfrak{gl}(n).
    \end{equation}
\end{proposition}

\begin{proof}
Let $X_1=X\otimes1$ and $X_2=1\otimes X$, and write
\begin{equation}
    X=\sum_{k,l}X_{kl}E_{kl}.
\end{equation}
By equation (\ref{eq:casimir_element}) we have
\begin{equation}
    [C_{12},X_1]
    =\sum_{i,j}[E_{ij},X]\otimes E_{ji} \quad \text{ and } \quad [C_{12},X_2]
    =\sum_{i,j}E_{ij}\otimes[E_{ji},X].
\end{equation}
Using
\begin{equation}\label{eq:[E_ij,E_kl]}
    [E_{ij},E_{kl}]
    =\delta_{jk}E_{il}-\delta_{il}E_{kj},
\end{equation}
we obtain
\begin{equation}
    [E_{ij},X]
    =\sum_l X_{jl}E_{il}-\sum_k X_{ki}E_{kj} \quad \text{ and } \quad 
    [E_{ji},X]
    =\sum_l X_{il}E_{jl}-\sum_k X_{kj}E_{ki}.
\end{equation}
Thus
\begin{align}
[C_{12},X_1]
    &=\sum_{i,j,l}X_{jl}E_{il}\otimes E_{ji}
    -\sum_{i,j,k}X_{ki}E_{kj}\otimes E_{ji}, \\
    [C_{12},X_2]
    &=\sum_{i,j,l}X_{il}E_{ij}\otimes E_{jl}
    -\sum_{i,j,k}X_{kj}E_{ij}\otimes E_{ki}.
\end{align}
After relabelling dummy indices, the first term equals the fourth and the second equals the third with opposite sign. Therefore
\begin{equation}
    [C_{12},X_1]+[C_{12},X_2]=0.
\end{equation}
\end{proof}

\begin{lemma}[Flip property of the Casimir tensor element]\label{lemma:flipping}
    Let
    \begin{equation}
        C_{12} = \sum_{i,j} E_{ij} \otimes E_{ji}
    \end{equation}
    be the Casimir tensor element of $\mathfrak{gl}(n)$.
    Then, for all $X,Y \in \mathfrak{gl}(n)$, one has
    \begin{equation}
        C_{12}(X \otimes Y)C_{12} = Y \otimes X.
    \end{equation}
\end{lemma}

\begin{proof}
    Write
    \begin{equation}
        X = \sum_{a,b} X_{ab} E_{ab},
        \qquad
        Y = \sum_{c,d} Y_{cd} E_{cd}.
    \end{equation}
    Then
    \begin{equation}
        C_{12}(X \otimes Y)C_{12}
        =
        \sum_{i,j,k,l}
        (E_{ij} X E_{kl}) \otimes (E_{ji} Y E_{lk}).
    \end{equation}
    Using
    \begin{equation}
        E_{ij} E_{ab} E_{kl}
        =
        \delta_{ja}\delta_{bk} E_{il},
        \qquad
        E_{ji} E_{cd} E_{lk}
        =
        \delta_{ic}\delta_{dl} E_{jk},
    \end{equation}
    we obtain
    \begin{equation}
        E_{ij} X E_{kl}
        =
        \sum_{a,b} X_{ab}\delta_{ja}\delta_{bk} E_{il}
        =
        X_{jk} E_{il},
    \end{equation}
    and similarly
    \begin{equation}
        E_{ji} Y E_{lk}
        =
        \sum_{c,d} Y_{cd}\delta_{ic}\delta_{dl} E_{jk}
        =
        Y_{il} E_{jk}.
    \end{equation}
    Therefore
    \begin{equation}
        C_{12}(X \otimes Y)C_{12}
        =
        \sum_{i,j,k,l}
        X_{jk} Y_{il}\,
        E_{il} \otimes E_{jk}.
    \end{equation}
    Rearranging the sums, we get
    \begin{equation}
        C_{12}(X \otimes Y)C_{12}
        =
        \left(\sum_{i,l} Y_{il} E_{il}\right)
        \otimes
        \left(\sum_{j,k} X_{jk} E_{jk}\right)
        =
        Y \otimes X.
    \end{equation}
\end{proof}

The following identity will play a crucial role in the derivation of the Yang--Baxter equation in Section~\ref{sec:YB}.
\begin{proposition}\label{prop:casimir_r_relation}
    Let $C_{12}$ be the Casimir tensor element and let $r \in \mathfrak{gl}(n)^{\otimes 2}$. Then
    \begin{equation}\label{eq:casimir_YB}
        [C_{12},r_{13}] + [C_{12},r_{23}]
        + [C_{13},r_{21}] + [C_{13},r_{23}]
        + [C_{23},r_{21}] + [C_{23},r_{31}] = 0.
    \end{equation}
\end{proposition}
\begin{proof}
    We rewrite the left-hand side of \eqref{eq:casimir_YB} as
    \begin{equation}\label{eq:casimir_accorpati}
        [C_{12}, r_{13} + r_{23}]
        + [C_{13}, r_{21} + r_{23}]
        + [C_{23}, r_{21} + r_{31}].
    \end{equation}
    We consider the first term. By Lemma~\ref{lemma:flipping}, the operator $C_{12}$ acts as the transposition of the first two tensor factors. Therefore
    \begin{equation}
        C_{12} r_{13} C_{12} = r_{23},
        \qquad
        C_{12} r_{23} C_{12} = r_{13}.
    \end{equation}
    Hence
    \begin{equation}
        C_{12}(r_{13}+r_{23})C_{12} = r_{13}+r_{23}.
    \end{equation}
    Since $C_{12}^2 = 1$, then
    \begin{equation}
        C_{12}^{-1} (r_{13} + r_{23}) C_{12} = r_{13} + r_{23} \implies [C_{12}, r_{13}+r_{23}] = 0.
    \end{equation}
    The same argument applied cyclically gives
    \begin{equation}
        [C_{13}, r_{21}+r_{23}] = 0,
        \qquad
        [C_{23}, r_{21}+r_{31}] = 0.
    \end{equation}
    Summing these three identities, we obtain \eqref{eq:casimir_YB}.
\end{proof}